% !TEX root = ../main.tex

\chapter{Kesimpulan dan Saran}

\section{Kesimpulan}

Berdasarkan implementasi dan eksperimen pada proyek ini, dapat disimpulkan bahwa:

\begin{enumerate}
    \item Komponen \textit{forward propagation} CNN, SimpleRNN, dan LSTM berhasil diimplementasikan secara modular \textit{from scratch} menggunakan NumPy serta kompatibel dengan bobot Keras melalui \texttt{set\_weights()}.
    \item Pipeline klasifikasi citra dan image captioning berhasil dibangun dari tahap preprocessing, pelatihan Keras, pemuatan bobot ke model scratch, hingga evaluasi metrik pada data uji.
    \item Evaluasi CNN menegaskan pentingnya desain arsitektur (jumlah layer, jumlah filter, ukuran kernel, dan jenis pooling) serta menunjukkan trade-off antara performa dan efisiensi parameter pada shared vs non-shared convolution.
    \item Evaluasi captioning menunjukkan pengaruh signifikan jumlah layer recurrent, ukuran hidden state, dan strategi decoding terhadap kualitas caption, dengan LSTM umumnya lebih stabil untuk dependensi urutan yang lebih panjang.
    \item Implementasi bonus (batch inference, backward propagation, init-inject, beam search, dan Grad-CAM) memperluas cakupan analisis dan meningkatkan kualitas interpretasi model.
\end{enumerate}

\section{Saran}

Saran pengembangan untuk pekerjaan selanjutnya adalah sebagai berikut.

\begin{enumerate}
    \item Menambahkan protokol evaluasi yang lebih kuat (misalnya \textit{multiple seeds} dan analisis statistik) agar kesimpulan performa lebih robust.
    \item Menjalankan \textit{error analysis} terstruktur pada sampel prediksi gagal, baik untuk klasifikasi maupun captioning, agar perbaikan arsitektur lebih terarah.
    \item Mengoptimasi performa komputasi melalui vektorisasi tambahan atau akselerasi perangkat keras saat inferensi batch besar.
    \item Mengembangkan decoder dengan teknik lanjutan (misalnya attention atau length normalization pada beam search) untuk meningkatkan kualitas caption.
    \item Memperluas dokumentasi eksperimen agar seluruh konfigurasi dan hasil dapat direproduksi dengan lebih mudah.
\end{enumerate}
